\setlength{\parindent}{0pt}

\chapter{Cartas, jaggerbombs y tortilleras}

\epigraf{La causa de tots els mals no és la ignorància, són els daddy issues.}{Regina Calsapeu}

\lettrine[loversize=0.2, lines=2]{C}{ristina} no solía revisar el buzón. Era uno de esos rituales de la vida adulta que siempre le daban pereza, como desatascar el fregadero o doblar las sábanas ajustables. Pero ese día algo la empujó. Tal vez una corazonada, tal vez el destino disfrazado de aburrimiento.\\

Entre facturas y propaganda del súper, había un sobre hecho a mano. Papel reciclado. Esquinas decoradas con flores prensadas. Y una caligrafía que le encogió el alma: Regina.\\

Su ex. Su historia inconclusa. Su terremoto emocional de hace tres años.\\

Cristina abrió la carta temblando, como quien abre una cápsula del tiempo que tal vez debería haberse quedado enterrada.\\

> “Cristina,\\

Sé que esto llega tarde, tal vez demasiado tarde. Pero no puedo callar más. Te pienso cada noche. Aún siento el eco de tus manos, de tu risa nerviosa cuando hablabas dormida.\\

Aún te amo.\\

No quiero interrumpir tu vida, solo saber si aún hay un espacio para mí en ella.\\

Tuya,\\
Regina.”\\


Cristina sintió un torbellino en el pecho. Había enterrado a Regina como quien entierra un meteorito: rápido, mal, y esperando que nunca más vuelva a emerger. Pero volvió.\\

\saltescena

Esa noche, mientras María dormía abrazando su copia de Stone Butch Blues, con el cabello desordenado y la respiración tranquila, Cristina no pudo dejar de pensar en Regina. En su caos, en su intensidad, en lo mal que acabaron... y lo bien que se sentía cuando estaban juntas.\\

Amaba a María. Pero el deseo, esa parte salvaje del amor, tenía su propio idioma.\\

Con el corazón roto y temblor en los dedos, Cristina recogió sus cosas: una mochila con ropa, Ácido el pez —que le hacía burbujitas como diciéndole “¿ya otra vez, amiga?”—, y el libro de María, el favorito, para no olvidarla nunca.\\

Dejó su anillo en la mesita. Y una carta.\\

> “María,\\

Perdóname. No sé cómo explicar esto sin que duela, pero la carta que recibí cambió todo. Regina volvió a mi vida. Y no puedo negarme a escuchar lo que me grita el pecho.\\

Me duele dejarte. Eres lo más estable y hermoso que he tenido. Pero hay amores que queman desde dentro y uno no elige si quiere apagar el incendio.\\

Me llevo algo tuyo, y no hablo solo de Ácido, sino de todo lo que me enseñaste.\\

Con cariño eterno,\\
C.”\\

 

Salió de casa antes del amanecer. Llamó a David.\\

—¿Me llevas al aeropuerto?\\
—Voy con el depósito lleno. ¿Qué ha pasado?\\

Cristina no contestó. Solo abrazó fuerte su mochila cuando vio los faros del coche aparecer.\\

Pero justo cuando estaba por subirse, un coche estacionado chirrió. De él salió Jordi, su ex de hace años, con energía de “yo hago meditación pero sigo siendo un intenso”.\\

—¡Cristina! ¡No hagas esto! ¡Te estás arrepintiendo antes de subir al coche!\\
—No, Jordi. Ya me arrepentí demasiado en silencio. Ahora actúo.\\

David, desde el coche, puso Miley Cyrus – “The Climb” para hacer más épica la escena.\\

Y se fueron.\\

\saltescena

María despertó con una sensación extraña. La cama fría. El silencio. Y la ausencia.\\

Al ver la carta y el anillo, no lloró. No gritó. Se quedó quieta, como si su cuerpo necesitara reiniciar. Cuando finalmente reaccionó, se vistió con lo primero que encontró: su chaqueta de cuero, su pantalón negro lleno de pelusa, y se fue al bar más cercano.\\

El Armario Abierto era un refugio para corazones rotos y bolleras poetas. María pidió un jaggerbomb. Luego otro. Luego uno más, con Red Bull y lágrimas.\\

Estaba por rendirse a la tristeza cuando vio una cara conocida.\\

—¿Pollalex? ¿Lhuevos?\\
—¡María! ¿Qué te ha pasado?\\
—Cristina me dejó por una ex. Se llevó a Ácido. Y mi puto libro.\\
—Uy. Eso amerita mínimo cuatro tragos y una ronda de insultos creativos —dijo Lhuevos, que era especialista en eso.\\

La conversación apenas comenzaba cuando entró Vicente, el odiado profesor de fisiología. Ex-novio de la heteronorma. Enemigo de la sororidad.\\

—¡¿Qué hacéis aquí, tortilleras?! ¡A la cocina! ¡Esto es degeneración pura!\\

María se levantó. Estaba triste, sí, pero jamás lo suficiente como para dejar pasar eso.\\

—Ven aquí, pedazo de célula muerta —dijo, y se le tiró encima.\\

La pelea fue legendaria. Hubo puñetazos, gritos, una silla volando, y alguien gritó “¡por Sappho!”.\\
La policía vegana llegó. Eran estrictas, pacíficas y tenían tofu en vez de porras.\\

—Tú, la de la chaqueta de cuero. ¿Eres tortillera?\\
—Sí, ¿y qué?\\
—Entonces has comido tortilla.\\
—¡No he comido tortilla! ¡Soy intolerante al huevo!\\
—Demasiado específico. Detenida.\\

\saltescena

Horas después, María estaba en una celda, con el maquillaje corrido y la moral por los suelos. Pensó en Ácido. En Stone Butch Blues. En cómo todo se había ido a la mierda tan rápido.\\

—¿Puedo saber quién pagó mi fianza? —preguntó, mientras le devolvían sus cosas.\\

La puerta se abrió.\\
Y ahí estaba.\\

Marta Ruiz.\\
Su primer amor. Su error favorito. Su poema inconcluso.\\
Con ese abrigo azul eléctrico, sus ojos con brillo de fuego suave y ese pelo que parecía peinado por musas.\\

—Hola, María —dijo, con voz grave y cálida.\\
—¿Tú… tú qué haces aquí?\\
—Escuché lo que pasó. Y no iba a dejar que salieras sola del agujero.\\

María sintió que algo viejo y bonito despertaba dentro.\\
—Sigo sin comer tortilla, ¿eh?\\
—Perfecto. Yo sigo sin cocinarla.\\

Y se abrazaron. Lento. Real. Como quien ha perdido mucho pero aún cree en las segundas, terceras y cuartas oportunidades.

